\chapter{Einführung}
\label{chap:einf}
Software-Systeme sind ein essenzieller Bestandteil der Wirtschaft und besitzen hohe Relevanz für Endkunden. Bereits ein lokaler Ausfall in einem komplexen System kann katastrophale Folgen für Firmen und Privatpersonen haben. Ein Ausfall im \ac{DNS}-System des Cloudanbieters \ac{AWS} \cite{Staff_2025} zeigt, dass schon geringfügige Störungen zu globalen \ac{IT}-Problemen führen. Da viele IT-Produkte auf Cloudlösungen basieren, verursachen Ausfälle dieser Systeme oft signifikante Schäden.

Die in den letzten Jahren entstandene DevOps-Bewegung \cite{bass2015devops} zielt darauf ab, die Softwarestabilität in der Produktionsphase zu gewährleisten und zeitnahe Reaktionen auf Fehler zu ermöglichen. Entwicklung und Betrieb von Software gelten nicht mehr als getrennte Bereiche, da ihre Symbiose zunehmend an Bedeutung gewinnt.



\section{Motivation: Herausforderungen der unstrukturierten Log-Analyse}
System-Logs spielen dabei eine wichtige Rolle, da Inhalt und Häufigkeit den Softwarezustand abbilden. Dabei unterscheidet man zwischen Logs des unterliegenden Systems (Server) und Nachrichten der Software selbst (application logs). Linux-System-Logs folgen oft dem klaren Format des Syslog-Protokolls \cite{rfc5424}. Das erleichtert die Verarbeitung, da diese Syslogs einem durch das Protokoll definierten Standard entsprechen.
Individuelle Lognachrichten aus eigener Softwareentwicklung folgen hingegen oft keinem festen Format, was Weiterverarbeitung und Analyse erschwert. Zur Verwaltung und Analyse großer Log-Mengen eignet sich Open-Source-Software wie der \ac{ELK}-Stack. Dieser besteht aus Elasticsearch (einer NoSQL-Datenbank für \ac{JSON}-Daten), Logstash (einer Pipeline zur Log-Weiterleitung) und Kibana (einer Weboberfläche zur Datenvisualisierung).
Kibana eignet sich gut für die Suche nach Log-Nachrichten, löst jedoch ein zentrales Problem nicht: Viele Logs unterscheiden sich nur durch Metadaten wie Zeitstempel oder \ac{IP}-Adressen, besitzen aber die gleiche Aussage.

\begin{logbox}{Beispiel 1: \ac{SSH} Login-Fehler - Root}
\begin{lstlisting}
Nov  6 06:12:05 servername sshd[12345]: Failed password for invalid user root from 191.167.10.80 port 4022 ssh2
\end{lstlisting}
\end{logbox}

\begin{logbox}{Beispiel 2: \ac{SSH} Login-Fehler - Admin}
\begin{lstlisting}
Nov  6 09:11:05 servername sshd[12345]: Failed password for invalid user admin from 192.168.1.100 port 22 ssh2
\end{lstlisting}
\end{logbox}


Die Beispiele veranschaulichen Log-Nachrichten, die ähnliche Ereignisse beschreiben. Eine Gruppierung solcher Logs ist sinnvoll. Dies gelingt beispielsweise mit regulären Ausdrücken (\acs{Regex}), die variable Teile filtern und Gruppen bilden. Dieses Verfahren bietet Flexibilität bei der individuellen Suche, erfordert jedoch hohen manuellen Pflegeaufwand, etwa bei neuen Log-Typen. Alternativ gruppieren Algorithmen den Text der Log-Nachrichten automatisch. Diese Arbeit untersucht Möglichkeiten, Log-Nachrichten kontinuierlich und mit minimalem manuellen Aufwand zu clustern.

\section{Zielsetzung: Vergleich von Verfahren für kontinuierliches Log-Clustering} \label{forschungs_frage}
Ziel dieser Arbeit ist die Beantwortung der Frage, welches Verfahren sich für die automatisierte, kontinuierliche Gruppierung von Lognachrichten am besten eignet - hinsichtlich Cluster-Qualität, Reihenfolge-Stabilität und
struktureller Eignung für den kontinuierlichen Betrieb. Dazu werden sechs Ansätze -- vier algorithmische Verfahren (K-Means, DBSCAN, Drain3, Brain), ein probabilistischer Ansatz (MinHash mit Locality-Sensitive Hashing) sowie ein vektorbasierter Ansatz auf Basis von Sentence-Embeddings und PostgreSQL mit \texttt{pgvector}-Erweiterung verglichen.

\section{Auswahl der untersuchten Verfahren}
\label{sec:auswahl}
Die untersuchten Verfahren wurden entlang der für das Log-Clustering 
relevanten methodischen Ansätze ausgewählt, um eine möglichst breite 
Abdeckung des Lösungsraums zu gewährleisten. Berücksichtigt wurden 
sowohl klassische Verfahren des unüberwachten Lernens als auch 
domänenspezifische Logparser sowie ein vektorbasierter Ansatz auf 
Grundlage neuronaler Sprachmodelle.

Aus dem Bereich des klassischen Clusterings repräsentiert K-Means
(Abschnitt~\ref{sec:kmeans}) das partitionierende vorgehen. DBSCAN
(Abschnitt~\ref{sec:dbscan}) ergänzt diese Perspektive um ein
dichtebasiertes Verfahren, das ohne die Vorgabe einer Cluster-Anzahl
auskommt.

Aus dem Bereich der domänenspezifischen Logparser wurden Drain3 
(Abschnitt~\ref{sec:drain3}) und Brain (Abschnitt~\ref{sec:brain}) 
ausgewählt. Drain3 gilt in der Literatur als etablierter Online-Parser 
mit fester Baumtiefe, während Brain eine konzeptionelle Schwäche von 
Drain -- die Abhängigkeit von der Token-Position -- explizit adressiert. 

MinHash mit Locality-Sensitive Hashing repräsentiert das probabilistische 
Paradigma und ermöglicht sublineare Ähnlichkeitssuche auf großen 
Datenmengen. Der vektorbasierte Ansatz auf Basis von 
Sentence-Embeddings und der PostgreSQL-Erweiterung \texttt{pgvector} 
bildet schließlich den Übergang zu modernen, semantisch arbeitenden 
Verfahren ab (Kapitel~\ref{chap:vek}).

\section{Methodik: Evaluierung mittels ARI-Wert und V-Measure} \label{Vorgensweise}

Die Qualität der Clustering-Algorithmen wird durch zwei Metriken bestimmt: der \ac{ARI}-Wert \cite{gates2017impactrandommodelsclustering} und das V-Measure \cite{inproceedings}.

\paragraph{Adjusted Rand Index.}
Der \acs{ARI} \cite{YANG201719} misst, wie gut zwei
Partitionierungen -- die vorhergesagte Clusterzugehörigkeit und die
Ground Truth -- auf Ebene von Datenpunktpaaren übereinstimmen. Dazu
wird berechnet, welcher Anteil aller Paare in beiden Partitionierungen
übereinstimmend in denselben oder in verschiedene Cluster eingeteilt
ist. Als Weiterentwicklung des ursprünglichen Rand-Index (\acs{RI}) korrigiert
der ARI den Zufallseffekt statistisch: Der gemessene Rand-Index wird
mit dem Erwartungswert bei zufälliger Verteilung verglichen und um
diesen normalisiert. Diese Korrektur ist notwendig, weil der
ursprüngliche Rand-Index seinen Wertebereich ineffizient nutzt --
selbst zufällige Clusterings liefern Werte nahe~$1$, was die
Interpretation erschwert. Durch die Normalisierung erlaubt der ARI eine
klare Interpretation: Ein Wert von $1{,}0$ bedeutet vollständige
Übereinstimmung mit der Ground Truth, ein Wert nahe~$0$ entspricht
einer zufälligen Zuordnung, negative Werte zeigen eine schlechtere
Übereinstimmung als der Zufall an.

\paragraph{V-Measure.}
Das V-Measure ist eine entropiebasierte Metrik zur externen Evaluierung von Cluster-Analysen \cite{inproceedings}. „Extern“ bedeutet, dass das Ergebnis mit einer bekannten, korrekten Klasseneinteilung (Ground Truth) verglichen wird. Das V-Measure behebt Schwächen anderer Metriken, insbesondere die Abhängigkeit von Cluster-Anzahl oder Datensatzgröße. Es bildet das harmonische Mittel aus Homogenität und Vollständigkeit: Homogenität liegt vor, wenn jeder Cluster ausschließlich Datenpunkte derselben Klasse enthält, Vollständigkeit hingegen, wenn alle Datenpunkte einer Klasse demselben Cluster zugeordnet sind \cite[S.~410]{inproceedings}.
Die Versuche nutzen die V-Measure-Werte der Python-Bibliothek „scikit-learn“, die wie folgt berechnet werden: \footnote{Dokumentation der Funktion v\_measure\_score \url{https://scikit-learn.org/stable/modules/generated/sklearn.metrics.v_measure_score.html}}
\begin{equation*}
  v = \frac{(1 + \beta) \cdot \text{homogeneity} \cdot \text{completeness}}{\beta \cdot \text{homogeneity} + \text{completeness}}
\end{equation*}
Der Beta-Wert ist ein variabler Faktor (Standardwert 1,0), der die Gewichtung von Homogenität und Vollständigkeit anpasst.
Das V-Measure beurteilt Homogenität und Vollständigkeit und analysiert, ob Klassen vermischt oder zersplittert sind. Der ARI korrigiert hingegen den Zufallseinfluss und dient als strenger Benchmark für die Übereinstimmung mit der Referenz.



\section{Abgrenzung}
Die Arbeit entwirft keine neuen Algorithmen, sondern nutzt bestehende, etablierte und getestete Bibliotheken. Im Vordergrund steht der methodische Vergleich der Verfahren hinsichtlich der in Abschnitt~\ref{Vorgensweise} definierten Qualitäts- und Stabilitätskriterien. Quantitative Messungen von Laufzeit und Ressourcenverbrauch werden bewusst ausgeklammert, da eine reproduzierbare Performance-Bewertung eine isolierte Hardware-Umgebung sowie ein dediziertes Benchmark-Setup voraussetzt, deren methodischer Aufwand den Rahmen dieser Arbeit überschreitet. Stattdessen werden Aspekte des Ressourcenverhaltens qualitativ über die strukturelle Eignung für den inkrementellen Betrieb diskutiert. Ebenso ist die Entwicklung eines produktiven Endsystems mit Web-Oberfläche oder Anbindung an eine Log-Aggregations-Pipeline (etwa Elasticsearch) nicht Gegenstand dieser Arbeit; die Implementierung erfolgt ausschließlich in einer Jupyter-Notebook-basierten Evaluationsumgebung.

%\section{Struktur der Arbeit}

